% GENERATED FILE. Edit canonical lessons, then run npm run generate:books.
% canonical-source-hash: fnv1a64:12a04c2e047e58f5
% canonical-lessons: TE-C270-tapela, TE-C270-kundi, TE-C270-penam, TE-C270-rolu, TE-C270-rokali

\chapter{Saucepan, Flowerpot, Griddle, Mortar, Pestle}
\label{ch:te-saucepan-flowerpot-griddle-mortar-pestle}
\hlchaptermodality{270}

\begin{chapteropening}
\textbf{By the end of this chapter:} \emph{I can say a saucepan, a flowerpot, a griddle, a mortar and a pestle.}

Complete the last lesson of chapter 270: I can say a saucepan, a flowerpot, a griddle, a mortar and a pestle.
\end{chapteropening}

\section[\texorpdfstring{\te{తపేలా} (tapēlā)}{తపేలా (tapēlā)}]{\te{తపేలా} (tapēlā) — a saucepan}

\label{lesson:TE-C270-tapela}



\begin{quote}
Before the new one: say the Telugu for wood, then the Telugu for brass.
\end{quote}

\subsection*{You'll want to know: \te{తపేలా}}
\textbf{\te{తపేలా}} — \emph{tapēlā} — \textquotedblleft{}a saucepan\textquotedblright{}.

\begin{grammarlens}[title={What you've built}]
One more word for this chapter.
\end{grammarlens}

\subsection*{Your turn}
\begin{itemize}
\raggedright
  \item \emph{Say it:} \emph{tapēlā}
  \item \emph{Say it:} \emph{tapēlā}, once more
  \item \emph{Say it:} say \emph{tapēlā}
  \item \emph{Recall:} say the Telugu for wood, then the Telugu for brass, then say \emph{tapēlā} again
\end{itemize}

\begin{tcolorbox}[breakable,colback=teal!4,colframe=teal!35!black,title={Before you move on}]
What does \te{తపేలా} mean? (\textquotedblleft{}A saucepan\textquotedblright{}.) Say it once more.
\end{tcolorbox}
\section[\texorpdfstring{\te{కుండీ} (kuṇḍī)}{కుండీ (kuṇḍī)}]{\te{కుండీ} (kuṇḍī) — a flowerpot}

\label{lesson:TE-C270-kundi}



\begin{quote}
Before the new one: say the Telugu for brass, then the Telugu for a saucepan.
\end{quote}

\subsection*{You'll want to know: \te{కుండీ}}
\textbf{\te{కుండీ}} — \emph{kuṇḍī} — \textquotedblleft{}a flowerpot\textquotedblright{}.

\begin{grammarlens}[title={What you've built}]
One more word for this chapter.
\end{grammarlens}

\subsection*{Your turn}
\begin{itemize}
\raggedright
  \item \emph{Say it:} \emph{kuṇḍī}
  \item \emph{Say it:} \emph{kuṇḍī}, once more
  \item \emph{Say it:} say \emph{kuṇḍī}
  \item \emph{Recall:} say the Telugu for brass, then the Telugu for a saucepan, then say \emph{kuṇḍī} again
\end{itemize}

\begin{tcolorbox}[breakable,colback=teal!4,colframe=teal!35!black,title={Before you move on}]
What does \te{కుండీ} mean? (\textquotedblleft{}A flowerpot\textquotedblright{}.) Say it once more.
\end{tcolorbox}
\section[\texorpdfstring{\te{పెనం} (penaṁ)}{పెనం (penaṁ)}]{\te{పెనం} (penaṁ) — a griddle (tava)}

\label{lesson:TE-C270-penam}



\begin{quote}
Before the new one: say the Telugu for a saucepan, then the Telugu for a flowerpot.
\end{quote}

\subsection*{You'll want to know: \te{పెనం}}
\textbf{\te{పెనం}} — \emph{penaṁ} — \textquotedblleft{}a griddle (tava)\textquotedblright{}.

\begin{grammarlens}[title={What you've built}]
One more word for this chapter.
\end{grammarlens}

\subsection*{Your turn}
\begin{itemize}
\raggedright
  \item \emph{Say it:} \emph{penaṁ}
  \item \emph{Say it:} \emph{penaṁ}, once more
  \item \emph{Say it:} say \emph{penaṁ}
  \item \emph{Recall:} say the Telugu for a saucepan, then the Telugu for a flowerpot, then say \emph{penaṁ} again
\end{itemize}

\begin{tcolorbox}[breakable,colback=teal!4,colframe=teal!35!black,title={Before you move on}]
What does \te{పెనం} mean? (\textquotedblleft{}A griddle (tava)\textquotedblright{}.) Say it once more.
\end{tcolorbox}
\section[\texorpdfstring{\te{రోలు} (rōlu)}{రోలు (rōlu)}]{\te{రోలు} (rōlu) — a mortar (for pounding)}

\label{lesson:TE-C270-rolu}



\begin{quote}
Before the new one: say the Telugu for a flowerpot, then the Telugu for a griddle.
\end{quote}

\subsection*{You'll want to know: \te{రోలు}}
\textbf{\te{రోలు}} — \emph{rōlu} — \textquotedblleft{}a mortar (for pounding)\textquotedblright{}.

\begin{grammarlens}[title={What you've built}]
One more word for this chapter.
\end{grammarlens}

\subsection*{Your turn}
\begin{itemize}
\raggedright
  \item \emph{Say it:} \emph{rōlu}
  \item \emph{Say it:} \emph{rōlu}, once more
  \item \emph{Say it:} say \emph{rōlu}
  \item \emph{Recall:} say the Telugu for a flowerpot, then the Telugu for a griddle, then say \emph{rōlu} again
\end{itemize}

\begin{tcolorbox}[breakable,colback=teal!4,colframe=teal!35!black,title={Before you move on}]
What does \te{రోలు} mean? (\textquotedblleft{}A mortar (for pounding)\textquotedblright{}.) Say it once more.
\end{tcolorbox}
\section[\texorpdfstring{\te{రోకలి} (rōkali)}{రోకలి (rōkali)}]{\te{రోకలి} (rōkali) — a pestle}

\label{lesson:TE-C270-rokali}



\begin{quote}
Before the new one: say the Telugu for a griddle, then the Telugu for a mortar.
\end{quote}

\subsection*{You'll want to know: \te{రోకలి}}
\textbf{\te{రోకలి}} — \emph{rōkali} — \textquotedblleft{}a pestle\textquotedblright{}.

\begin{grammarlens}[title={What you've built}]
One more word for this chapter.
\end{grammarlens}

\subsection*{Your turn}
\begin{itemize}
\raggedright
  \item \emph{Say it:} \emph{rōkali}
  \item \emph{Say it:} \emph{rōkali}, once more
  \item \emph{Say it:} say \emph{rōkali}
  \item \emph{Recall:} say the Telugu for a griddle, then the Telugu for a mortar, then say \emph{rōkali} again
\end{itemize}

\begin{tcolorbox}[breakable,colback=teal!4,colframe=teal!35!black,title={Before you move on}]
What does \te{రోకలి} mean? (\textquotedblleft{}A pestle\textquotedblright{}.) Say it once more.
\end{tcolorbox}
